\section{Level I: Dual Numbers}

The resulting framework exposes a convenient high-level interface
through the Python programming language which can be used to state,
compose, and efficiently deploy complex physical models. 
Not only does
this open the door for powerful new computational procedures and inverse
analyses, but it substantially simplifies how we interact with FE
software, and provides an accessible tool for engineering education.

% \hypertarget{part-i-composition-of-expressions}{\chapter{Part I: Composition of Expressions}\label{part-i-composition-of-expressions}}

In \textcite{frostig2018compiling}, researchers from the Google Brain
team begin their presentation of the JAX project with a
simple observation; machine learning workloads tend to consist largely
of subroutines satisfying two simple properties:

\begin{enumerate}
\def\labelenumi{\arabic{enumi}.}
\tightlist
\item
  They are \emph{pure} in the sense that their execution is free of side
  effects, and
\item
  They are \emph{statically composed} in the sense that they implement a
  computation which can be defined as a static data dependency on a
  graph of some set of \emph{primitive} operations (This will be further
  explained).
\end{enumerate}

The exploitation of these properties with techniques such as automatic
differentiation and just-in-time compilation has become a central
component of tools like TensorFlow, Matlab and PyTorch, which in the
recent decade have been at the forefront of an ``explosion'' in the
field of deep learning. The utility that these ground-breaking platforms
provide is that they present researchers and developers with a high
level platform for concisely staging mathematical computations with an
intuitive ``neural network'' paradigm. By their structure, these
resulting models admit analytic gradients which allows for the
application of powerful optimization algorithms.


\subsection{Dual Numbers}\label{building-a-differentiable-model}

There are several approaches which may be followed to arrive at the same
formulation of forward mode automatic differentiation
\autocite{hoffmann2016hitchhiker}. A particularly elegant entry point
for the purpose of illustration is that used by
\textcite{pearlmutter2007lazy} which begins by considering the algebra
of \emph{dual numbers}. 
These numbers are expressions of the form
\(a + b\varepsilon\), with \(a,b \in \mathbb{R}\) for which addition
and multiplication is like that of complex numbers, save that the symbol
\(\varepsilon\) be defined to satisfy \(\varepsilon^2 = 0\).

Consider an arbitrary analytic function
\(f: \mathbb{R} \rightarrow \mathbb{R}\), admitting the following Taylor series expansion:
%
\[
f(a+b \varepsilon)=\sum_{n=0}^{\infty} \frac{f^{(n)}(a) b^{n} \varepsilon^{n}}{n !} 
\]
%
By truncating the series at \(n = 1\) and adhering to the operations
defined over dual number arithmetic, one obtains the expression
%
\[
f(a + b \varepsilon) = f(a) + bf^\prime (a) \varepsilon, 
\]
%
This result insinuates that by elevating such a function to perform on
dual numbers, it has effectively been transformed into a new function
which simultaneously evaluates both it's original value, and it's
derivative \(f^\prime\).

For instance, consider the following pair of mappings:

\[
\begin{aligned}
f: & \mathbb{R} \rightarrow \mathbb{R} \\
   & x \mapsto (c - x)
\end{aligned}
\]

\begin{equation}
\begin{aligned}
g: & \mathbb{R} \rightarrow \mathbb{R} \\
   & x \mapsto a x + b f(x)
\end{aligned}
\label{eq:dual-area}\end{equation}
%
One can use the system of dual numbers to evaluate the gradient of \(g\)
at some point \(x \in \mathbb{R}\), by evaluating \(g\) at the dual
number \(x + \varepsilon\), where it is noted that the dual coefficient
\(b = 1\) is in fact equal to the derivative of the dependent variable
\(x\). This procedure may be thought of as \emph{propagating} a
differential \(\varepsilon\) throughout the computation.

\[
\begin{aligned}
g(x + \varepsilon)   & = a (x + \varepsilon) + b f(x + \varepsilon ) \\
                     & = a x + a \varepsilon + b ( c - x - \varepsilon) \\
                     & = (a - b) x + bc + (a - b) \varepsilon 
\end{aligned}
\]

The result of this expansion is a dual number whose real component is
the value of \(g(x)\), and dual component, \(a-b\), equal to the
gradient of \(g\) evaluated at \(x\).

\begin{figure}
\hypertarget{fig:area-graph}{%
\centering
\includegraphics{Figures/area.pdf}
\caption{Computational graph for the area of a
T-section.}\label{fig:area-graph}
}
\end{figure}

Such a system lends itself well to implementation in programming
languages where an object might be defined to transparently act as a
scalar. For example, in C, one might create a simple data structure with
fields \texttt{real} and \texttt{dual} that store the real component,
\(a\) and \emph{dual} component, \(b\) of such a number as show in \cref{lst:struct-dual}.

\hypertarget{lst:struct-dual}{%
\label{lst:struct-dual}}%
\begin{Shaded}
\begin{Highlighting}[]
\KeywordTok{typedef} \KeywordTok{struct}\NormalTok{ dual }\OperatorTok{\{}
    \DataTypeTok{float}\NormalTok{ real}\OperatorTok{,}\NormalTok{ dual}\OperatorTok{;}
\OperatorTok{\}} \DataTypeTok{dual\_t}\OperatorTok{;}
\end{Highlighting}
\end{Shaded}

Then, using the operator overloading capabilities of C++, the
application of the \texttt{*} operator can be overwritten as follows so
that the property \(\varepsilon = 0\) is enforced, as show in
\cref{lst:dual-mult}.

\hypertarget{lst:dual-mult}{%
\label{lst:dual-mult}}%
\begin{Shaded}
\begin{Highlighting}[]
\DataTypeTok{dual\_t} \DataTypeTok{dual\_t}\OperatorTok{::}\KeywordTok{operator} \OperatorTok{*} \OperatorTok{(}\DataTypeTok{dual\_t}\NormalTok{ a}\OperatorTok{,} \DataTypeTok{dual\_t}\NormalTok{ b}\OperatorTok{)\{}
    \ControlFlowTok{return} \OperatorTok{(}\DataTypeTok{dual\_t}\OperatorTok{)} \OperatorTok{\{}
        \OperatorTok{.}\NormalTok{real }\OperatorTok{=}\NormalTok{ a}\OperatorTok{.}\NormalTok{real }\OperatorTok{*}\NormalTok{ b}\OperatorTok{.}\NormalTok{real}\OperatorTok{,}
        \OperatorTok{.}\NormalTok{dual }\OperatorTok{=}\NormalTok{ a}\OperatorTok{.}\NormalTok{real }\OperatorTok{*}\NormalTok{ b}\OperatorTok{.}\NormalTok{dual }\OperatorTok{+}\NormalTok{ b}\OperatorTok{.}\NormalTok{real }\OperatorTok{*}\NormalTok{ a}\OperatorTok{.}\NormalTok{dual}
    \OperatorTok{\};}
\OperatorTok{\}}
\end{Highlighting}
\end{Shaded}

Similar operations are defined for variations of the \texttt{*},
\texttt{+}, and \texttt{-} operators in a single-file library in the
appendix. With this complete arithmetic interface implemented, forward
propagation may be performed simply with statements like that of
\cref{lst:dual-expr}.

\hypertarget{lst:dual-expr}{%
\label{lst:dual-expr}}%
\begin{Shaded}
\begin{Highlighting}[]
\DataTypeTok{int}\NormalTok{ main}\OperatorTok{(}\DataTypeTok{int}\NormalTok{ argc}\OperatorTok{,} \DataTypeTok{char} \OperatorTok{**}\NormalTok{argv}\OperatorTok{)\{}
  \DataTypeTok{dual\_t}\NormalTok{ x}\OperatorTok{(}\FloatTok{6.0}\OperatorTok{,}\FloatTok{1.0}\OperatorTok{);}
  \DataTypeTok{dual\_t}\NormalTok{ dg }\OperatorTok{=}\NormalTok{ a }\OperatorTok{*}\NormalTok{ x }\OperatorTok{+} \OperatorTok{(}\NormalTok{c }\OperatorTok{{-}}\NormalTok{ x}\OperatorTok{)} \OperatorTok{*}\NormalTok{ b}\OperatorTok{;}
\NormalTok{  printd}\OperatorTok{(}\NormalTok{dA}\OperatorTok{);}
\OperatorTok{\}}
\end{Highlighting}
\end{Shaded}

\hypertarget{template-expressions}{%
\subsection{Function Composition}\label{template-expressions}}

The simple \texttt{dual\_t} numeric type developed in the previous
section allows both an expression, and its gradient to be evaluated
simultaneously with only the minor penalty that each value stores an
associated \texttt{dual} value. However, if such a penalty was imposed
on every value carried in a modeling framework, both memory and
computation demands would scale very poorly. In order to eliminate this
burden, the concept of a \emph{template expression} is introduced. This
is again illustrated using the C++ programming language in order to
maintain transparency to the corresponding machine instructions, but as
expanded in the following part, this technique becomes extremely
powerful when implemented in a context where access to the template
expression in a processed form, such as an abstract syntax tree, can be
readily manipulated. Examples of this might be through a higher-level
language offering introspective capabilities (e.g., Python or most
notably Lisp), or should one be so brave, via compiler extensions (an
option which is becoming increasingly viable with recent developments
and investment in projects like LLVM \autocite{lattner2004llvm}).

In \ref{lst:dual-template}, the function \(g\) has been defined as
a \emph{template} on generic types \texttt{TA}, \texttt{TB}, \texttt{TC}
and \texttt{TX}. In basic use cases, like that shown in
\ref{lst:dual-template}, a modern C++ compiler will deduce the
types of the arguments passed in a call to \texttt{G} so that dual
operations will only be used when arguments of type
\VERB|\NormalTok{dual\_t}| are passed. Furthermore, such an
implementation will produce a derivative of \texttt{G} with respect to
whichever parameter is of type \VERB|\NormalTok{dual\_t}|. For example,
in the assignment to variable \texttt{dg} in
\ref{lst:dual-template}, a \VERB|\NormalTok{dual\_t}| value is
created from variable \texttt{x} within the call to \texttt{G}, which
automatically causes the compiler to create a function with signature
\VERB|\NormalTok{dual\_t }\OperatorTok{(*)(}\NormalTok{dual\_t}\OperatorTok{,}\NormalTok{ real\_t}\OperatorTok{,}\NormalTok{ real\_t}\OperatorTok{,}\NormalTok{ real\_t}\OperatorTok{)}|
which propagates derivatives in \texttt{x}.


\hypertarget{lst:dual-template}{%
\label{lst:dual-template}}%
\begin{Shaded}
\begin{Highlighting}[]
\KeywordTok{template}\OperatorTok{\textless{}}\KeywordTok{typename}\NormalTok{ TX}\OperatorTok{,}\KeywordTok{typename}\NormalTok{ TB}\OperatorTok{,} \KeywordTok{typename}\NormalTok{ TA}\OperatorTok{,} \KeywordTok{typename}\NormalTok{ TC}\OperatorTok{\textgreater{}}
\KeywordTok{auto}\NormalTok{ G}\OperatorTok{(}\NormalTok{TX x}\OperatorTok{,}\NormalTok{ TB b}\OperatorTok{,}\NormalTok{ TA a}\OperatorTok{,}\NormalTok{ TC c}\OperatorTok{)\{}
\NormalTok{  printd}\OperatorTok{(}\NormalTok{b}\OperatorTok{*}\NormalTok{x}\OperatorTok{);}
\NormalTok{  printd}\OperatorTok{(}\NormalTok{c }\OperatorTok{{-}}\NormalTok{ x}\OperatorTok{);}
  \ControlFlowTok{return}\NormalTok{ a }\OperatorTok{*}\NormalTok{ x }\OperatorTok{+} \OperatorTok{(}\NormalTok{c }\OperatorTok{{-}}\NormalTok{ x}\OperatorTok{)} \OperatorTok{*}\NormalTok{b}\OperatorTok{;}
\OperatorTok{\}}

\DataTypeTok{int}\NormalTok{ main}\OperatorTok{(}\DataTypeTok{int}\NormalTok{ argc}\OperatorTok{,} \DataTypeTok{char} \OperatorTok{**}\NormalTok{argv}\OperatorTok{)\{}
  \DataTypeTok{real\_t}\NormalTok{ a }\OperatorTok{=} \FloatTok{60.0}\OperatorTok{,}\NormalTok{ c}\OperatorTok{=}\FloatTok{18.0}\OperatorTok{,}\NormalTok{ b}\OperatorTok{=}\FloatTok{18.0}\OperatorTok{;}
  \DataTypeTok{real\_t}\NormalTok{ x }\OperatorTok{=} \FloatTok{6.0}\OperatorTok{;}
  \DataTypeTok{real\_t}\NormalTok{ g }\OperatorTok{=}\NormalTok{ G}\OperatorTok{(}\NormalTok{x}\OperatorTok{,}\NormalTok{b}\OperatorTok{,}\NormalTok{a}\OperatorTok{,}\NormalTok{c}\OperatorTok{);}
  \DataTypeTok{dual\_t}\NormalTok{ dg }\OperatorTok{=}\NormalTok{ G}\OperatorTok{(}\DataTypeTok{dual\_t}\OperatorTok{(}\NormalTok{x}\OperatorTok{,}\DecValTok{1}\OperatorTok{),}\NormalTok{b}\OperatorTok{,}\NormalTok{a}\OperatorTok{,}\NormalTok{c}\OperatorTok{);}
\NormalTok{  printd}\OperatorTok{(}\NormalTok{dg}\OperatorTok{);}
\OperatorTok{\}}
\end{Highlighting}
\end{Shaded}

In this minimal example, the set of operations for which we have
explicitly defined corresponding dual operations constitute a set of
\emph{primitive} operations which may be composed arbitrarily in
expressions which will implicitly propagate machine precision first
derivatives. These compositions may include arbitrary use of control
structures like \VERB|\ControlFlowTok{while}| /
\VERB|\ControlFlowTok{for}| loops, branching via
\VERB|\ControlFlowTok{if}| statements, and by composition of such
operations, numerical schemes involving indefinite iteration.



\hypertarget{part-ii-the-finite-element-method}{%
\subsection{Finite Element Models}\label{part-ii-the-finite-element-method}}



A particularly powerful dialect of FE methods for problems in resilience
is the composable approach developed through works like
\textcite{spacone1996fibre} where the classical beam element is
generalized as a composition of arbitrary section models, which may in
turn be comprised of several material models. Platforms like
{[}FEDEASLab{]} and {[}OpenSees{]} are designed from the ground up to
natively support such compositions \autocite{mckenna2010nonlinear}. This
study builds on this composable architecture by defining a single
protocol for model components which can be used to nest or compose
elements indefinitely.

In order to establish an analogous structure to that of ANNs which is
sufficiently general to model geometric and material nonlinearities, we
first assert that such a FE program will act as a mapping,
\(\Phi: \Theta\times\mathcal{S}\times\mathcal{U}\rightarrow\mathcal{S}\times\mathcal{U}\),
which walks a function,
\(f: \Theta\times\mathcal{S}\times\mathcal{U}\rightarrow\mathcal{S}\times\mathcal{R}\)
(e.g.~by applying the Newton-Raphson method), until arriving at a
solution, \((S^*,U^*)=\Phi(\Theta,S_0,U_0)\) such that
\(f(\Theta,S_0,U^*) = (S^*, \mathbf{0})\) in the region of \(U_0\). In
problems of mechanics, \(f\) may be thought of as a potential energy
gradient or equilibrium statement, and the elements of \(\mathcal{U}\)
typically represent a displacement vector, \(\mathcal{R}\) a normed
force residual, \(\mathcal{S}\) a set of path-dependent state variables,
and \(\Theta\) a set of parameters which an analyst may wish to optimize
(e.g.~the locations of nodes, member properties, parameters, etc.). The
function \(f\) is composed from a set of \(i\) local functions
\(g_i: \theta\times s\times u \rightarrow s\times r\) that model the
response and dissipation of elements over regions of the domain, which
in turn may be comprised of a substructure or material function, until
ultimately resolving through elementary operations,
\(h_j: \mathbb{R}^n\rightarrow\mathbb{R}^m\) (e.g.~Vector products,
trigonometric functions, etc.).

In a mathematical sense, for a given solution routine, all of the
information that is required to specify a computation which evaluates
\(\Phi\), and some gradient, \(D_\Theta \Phi\), of this model with
respect to locally well behaved parameters \(\Theta\) is often defined
entirely by the combination of (1) a set of element response functions
\(g_i\), and (2) a graph-like data structure, indicating their
connectivity (i.e., the static data dependency of
\textcite{frostig2018compiling}) . For such problems, within a
well-structured framework, the procedures of deriving analytic gradient
functions and handling low-level machine instructions can be lifted
entirely from the analyst through abstraction, who should only be
concerned with defining the trial functions, \(g_i\), and establishing
their connectivity.
